% Unit 084 terjemahan Bahasa Indonesia dari chapter6.tex baris 2313--2457.
% Sumber: Methods of Algebra, Volume 2, commit
% 9a5803ff77dd3257484cb177f851a73770a59dd3 (CC BY 4.0).
% stable-unit-id: o014.aljabr2.chapter6.profinite-cohomology-a
% Cakupan: Bagian 6.12, modul mulus, definisi kohomologi profinit,
% dan kompleks standar kontinu sampai Teorema 6.12 pertama.

% segment-id: o014.li2.chapter6.profinite-cohomology-a.h001
\section{Kohomologi grup profinit}\label{sec:profinite-cohomology}
\hypertarget{o014.aljabr2.chapter6.profinite-cohomology-a}{ }

% segment-id: o014.li2.chapter6.profinite-cohomology-a.p001
Bagian ini mengandaikan definisi dan hasil dasar tentang grup profinit dari
\S\ref{sec:profinite-groups}.

% segment-id: o014.li2.chapter6.profinite-cohomology-a.p002
Misalkan $G$ grup profinit dan $M$ modul-$G$. Tidak sulit memeriksa bahwa
syarat-syarat berikut ekuivalen:
\begin{enumerate}[(i)]
	\item $M=\bigcup_KM^K$, dengan $K$ berkisar atas subgrup normal buka dari
	$G$;
	\item untuk setiap $x\in M$, subgrup penstabil $\Stab_G(x)$ bersifat buka;
	\item pemetaan aksi $G\times M\to M$ kontinu jika $M$ diberi topologi
	diskret.
\end{enumerate}

% segment-id: o014.li2.chapter6.profinite-cohomology-a.d001
\begin{definition}
	\index{modul-$G$!mulus (smooth)}
	\index[sym1]{G-Mod-infty@$G\dcate{Mod}^\infty$}
	Misalkan $G$ grup profinit. Jika salah satu syarat di atas berlaku, modul-$G$
	$M$ disebut \emph{mulus}\footnote{Karena~(iii), banyak pustaka menyebutnya
	diskret.}. Semua modul-$G$ mulus membentuk subkategori penuh
	$G\dcate{Mod}^\infty$ dari $G\dcate{Mod}$.
\end{definition}

% segment-id: o014.li2.chapter6.profinite-cohomology-a.p003
Meskipun syarat pada $G$ sendiri bersifat topologis, karakterisasi~(i) atau
~(ii) bagi modul-$G$ mulus sepenuhnya bersifat aljabar: kemulusan merupakan sifat
modul-$G$, bukan struktur tambahan. Dari sini mudah dilihat bahwa
$G\dcate{Mod}^\infty$ adalah subkategori abelian dari $G\dcate{Mod}$. Operasi
dalam Definisi~\ref{def:Res-Infl} mempunyai analog alami dalam keadaan ini.

% segment-id: o014.li2.chapter6.profinite-cohomology-a.d002
\begin{definition-proposition}
	Misalkan $\varphi:H\to G$ homomorfisme kontinu antargrup profinit dan
$M$ modul-$G$ mulus. Maka modul-$H$ $\varphi^*M$ tetap mulus. Ini memberikan
	funktor
	\[ \varphi^*:G\dcate{Mod}^\infty\to H\dcate{Mod}^\infty. \]
\end{definition-proposition}
\begin{proof}
	Untuk setiap subgrup normal buka $K\lhd G$, prapetanya
	$\varphi^{-1}(K)\lhd H$ juga normal dan buka. Jika $x\in M^K$, maka
	$x\in(\varphi^*M)^{\varphi^{-1}K}$. Jadi
	\[ \varphi^*M=
	\bigcup_{L\lhd H:\,\text{buka}}(\varphi^*M)^L, \]
	yang membuktikan bahwa $\varphi^*M$ mulus.
\end{proof}

% segment-id: o014.li2.chapter6.profinite-cohomology-a.p004
Perhatikan bahwa funktor $\varphi^*$ selalu eksak. Dua kasus khusus berikut
akan digunakan:
\begin{description}
	\item[Restriksi] Jika $H$ subgrup tertutup dari $G$, ambil $\varphi$ sebagai
	homomorfisme penyertaan $H\hookrightarrow G$. Funktor terkait ditulis
	\[ \Res^G_H:G\dcate{Mod}^\infty\to H\dcate{Mod}^\infty. \]
	\item[Inflasi] Untuk subgrup normal tertutup $H\lhd G$, ambil $\varphi$
	sebagai homomorfisme hasil bagi $G\twoheadrightarrow G/H$. Funktor terkait
	ditulis
	\[ \mathrm{Infl}^G_{G/H}:G/H\dcate{Mod}^\infty
	\to G\dcate{Mod}^\infty. \]

	Grup $H$ dan $G/H$ yang disebutkan di atas otomatis profinit, sehingga
	definisi ini sah.
\end{description}
\index[sym1]{ResGH}
\index[sym1]{InflGH}

% segment-id: o014.li2.chapter6.profinite-cohomology-a.d003
\begin{definition-proposition}\label{def:smooth-part}
	Untuk setiap modul-$G$ $M$, definisikan
	$M^\infty:=\bigcup_KM^K$, dengan $K$ berkisar atas subgrup normal buka dari
	$G$. Ini merupakan modul-$G$ mulus, dan $M$ mulus jika dan hanya jika
	$M=M^\infty$.

	Untuk setiap modul-$G$ mulus $M'$, berlaku
	\[ \Hom_G(M',M)=\Hom_G(M',M^\infty). \]
	Dengan kata lain,
	$\text{penyertaan}:G\dcate{Mod}^\infty
	\leftrightarrows G\dcate{Mod}:(\cdot)^\infty$ merupakan pasangan adjoin.
\end{definition-proposition}
\begin{proof}
	Jelas.
\end{proof}

% segment-id: o014.li2.chapter6.profinite-cohomology-a.p005
Unsur-unsur $M^\infty$ juga disebut \emph{unsur mulus} dari $M$. Karena
$(\cdot)^\infty$ adalah adjoin kanan dari funktor eksak, funktor ini
mempertahankan objek injektif
(Proposisi~\ref{prop:adjoint-injective-projective}).

% segment-id: o014.li2.chapter6.profinite-cohomology-a.t001
\begin{corollary}
	Kategori abelian $G\dcate{Mod}^\infty$ memiliki cukup banyak objek
	injektif.
\end{corollary}
\begin{proof}
	Misalkan $M$ modul-$G$ mulus dan sertakan $M$ ke dalam modul-$G$ injektif
	$I$. Maka $M=M^\infty\hookrightarrow I^\infty$, sedangkan $I^\infty$
	merupakan modul-$G$ mulus injektif.
\end{proof}

% segment-id: o014.li2.chapter6.profinite-cohomology-a.r001
\begin{remark}
	Pada umumnya $G\dcate{Mod}^\infty$ tidak memiliki cukup banyak objek
	projektif. Inilah alasan utama bagian ini membahas kohomologi, bukan homologi.
\end{remark}

% segment-id: o014.li2.chapter6.profinite-cohomology-a.p006
Karena $G\dcate{Mod}^\infty$ memiliki cukup banyak objek injektif, kita dapat
membahas funktor turunan kanan dari funktor invarian $(\cdot)^G$. Pendekatan
bagian ini mula-mula memberikan definisi langsung yang mereduksi ke kasus grup
hingga, lalu mengidentifikasikannya dengan funktor turunan kanan dalam
Proposisi~\ref{prop:profinite-cohomology-identification}. Untuk sementara kita
hanya memakai bahasa klasik, tanpa kategori turunan.

% segment-id: o014.li2.chapter6.profinite-cohomology-a.p007
Mula-mula, ambil grup profinit $G$ dan subgrup normal buka $K$. Untuk setiap
modul-$G$ mulus $M$, submodul invarian-$K$
$M^K=(\Res^G_KM)^K$ secara alami menjadi modul-$G/K$. Dengan demikian
diperoleh funktor $(\cdot)^K:G\dcate{Mod}^\infty\to G/K\dcate{Mod}$ dan
pasangan adjoin yang jelas
\begin{equation}\label{eqn:invariant-adjunction-profinite}
	\begin{tikzcd}[row sep=tiny]
		\mathrm{Infl}^G_{G/K}:G/K\dcate{Mod} \arrow[shift left,r] &
		G\dcate{Mod}^\infty:(\cdot)^K \arrow[shift left,l]\\
		\Hom_G(\mathrm{Infl}^G_{G/K}(M'),M) \arrow[equal,r] &
		\Hom_{G/K}(M',M^K).
	\end{tikzcd}
\end{equation}

% segment-id: o014.li2.chapter6.profinite-cohomology-a.p008
Jika subgrup normal buka $K,K'$ dari $G$ memenuhi $K'\supset K$, maka
penyertaan $M^{K'}\hookrightarrow M^K$ ekuivarian terhadap homomorfisme hasil
bagi yang berarah sebaliknya $G/K'\twoheadleftarrow G/K$
(Definisi~\ref{def:group-equivariance}). Jadi untuk setiap
$n\in\Z_{\geq0}$ terdapat homomorfisme kanonik
\[ \rho^{K'}_K:\Hm^n(G/K',M^{K'})\to\Hm^n(G/K,M^K). \]
Jika $K_1\subset K_2\subset K_3$ merupakan subgrup normal buka dari $G$,
maka $\rho^{K_2}_{K_1}\rho^{K_3}_{K_2}=\rho^{K_3}_{K_1}$. Jadi definisi
berikut sah.

% segment-id: o014.li2.chapter6.profinite-cohomology-a.d004
\begin{definition}
	\index{kohomologi grup!kasus profinit}
	\index[sym1]{HnGM}
	Misalkan $G$ grup profinit dan $M$ modul-$G$ mulus. Untuk setiap
	$n\in\Z_{\geq0}$, definisikan
	\[ \Hm^n(G,M):=\varinjlim_K\Hm^n(G/K,M^K), \]
	dengan $K$ berkisar atas subgrup normal buka dari $G$, dan $\varinjlim$
	diambil terhadap urutan parsial
	$K'\preceq K\iff K\subset K'$ serta morfisme $\rho^{K'}_K$.
\end{definition}

% segment-id: o014.li2.chapter6.profinite-cohomology-a.p009
Semua subgrup normal buka membentuk himpunan terurut terfilter: himpunan ini
tertutup terhadap irisan hingga. Jika $G$ hingga, maka $\{1\}$ adalah satu-satunya
unsur maksimal, dan $\Hm^n(G,M)$ kembali menjadi versi awal dalam
\S\ref{sec:group-coh-sub}. Secara umum,
\[ \Hm^0(G,M)=\varinjlim_K(M^K)^{G/K}=M^G. \]

% segment-id: o014.li2.chapter6.profinite-cohomology-a.p010
Misalkan $f:M_1\to M_2$ homomorfisme modul-$G$ mulus. Untuk setiap subgrup
normal buka $K$, pemetaan ini membatasi menjadi
$f_K:M_1^K\to M_2^K$. Jelas bahwa $K'\supset K$ mengakibatkan
$\Hm^n(f_K)\rho^{K'}_K=\rho^{K'}_K\Hm^n(f_{K'})$, sehingga untuk setiap
$n\in\Z_{\geq0}$ terdapat
\[ \Hm^n(f):\Hm^n(G,M_1)\to\Hm^n(G,M_2). \]

% segment-id: o014.li2.chapter6.profinite-cohomology-a.e001
\begin{example}
	Misalkan $E|F$ perluasan Galois medan. Karena
	$E=\bigcup_{L|F}E^{\Gal(E|L)}=\bigcup_{L|F}L$, dengan $L|F$ menjelajahi
	subekstensi Galois hingga, aksi $\Gal(E|F)$ pada $E$ bersifat mulus.
	Dengan menerapkan Contoh~\ref{eg:additive-90-finite} pada setiap $L|F$,
	diperoleh
	\[ n\geq1\implies\Hm^n(\Gal(E|F),E)
	=\varinjlim_{L|F}\Hm^n(\Gal(L|F),L)=0. \]
\end{example}

% segment-id: o014.li2.chapter6.profinite-cohomology-a.p011
Seperti Proposisi~\ref{prop:groupcoh-cochain} dan
Catatan~\ref{rem:nor-cochain} untuk grup tanpa topologi, kohomologi grup
profinit dapat dihitung melalui kompleks standar atau versi ternormalisasinya;
perbedaannya terletak pada syarat kontinuitas.

% segment-id: o014.li2.chapter6.profinite-cohomology-a.d005
\begin{definition}
	\index{kompleks standar}
	\index[sym1]{CGM}
	Untuk modul-$G$ mulus $M$, beri $M$ topologi diskret dan definisikan
	\emph{kompleks standar} $C(G,M)=(C^n(G,M))_n$ dengan
	\[ C^n(G,M):=\{\text{pemetaan kontinu }f:G^n\to M\},
	\qquad n\in\Z_{\geq0}. \]
	Suku berderajat negatif didefinisikan sebagai $0$, sedangkan
	$d^n:C^n(G,M)\to C^{n+1}(G,M)$ diberikan oleh
	\begin{multline*}
		(d^nf)(g_1,\ldots,g_{n+1})
		=g_1f(g_2,\ldots,g_{n+1})\\
		+\sum_{k=1}^n(-1)^kf(\ldots,g_kg_{k+1},\ldots)
		+(-1)^{n+1}f(g_1,\ldots,g_n).
	\end{multline*}

	\index[sym1]{CGMbar}
	Definisikan \emph{kompleks standar ternormalisasi} sebagai subkompleks
	$\overline{C}(G,M)$ dari $C(G,M)$ berikut
	(lihat Catatan~\ref{rem:nor-cochain}):
	\[ \overline{C}^n(G,M)=\left\{\begin{array}{r|l}
		f\in C^n(G,M)&\exists,1\leq k\leq n,\ g_k=1_G\\
		&\implies f(g_1,\ldots,g_n)=0
	\end{array}\right\}. \]
\end{definition}

% segment-id: o014.li2.chapter6.profinite-cohomology-a.p012
Untuk setiap subgrup normal buka $K\lhd G$, terdapat penyertaan kompleks yang
jelas
$C(G/K,M^K)\hookrightarrow C(G,M)$ dan
$\overline{C}(G/K,M^K)\hookrightarrow\overline{C}(G,M)$.

% segment-id: o014.li2.chapter6.profinite-cohomology-a.p013
Karena $M$ diskret, pemetaan $f:G^n\to M$ kontinu jika dan hanya jika lokal
konstan. Kekompakan $G$ kemudian memberikan subgrup normal buka
$K_1\lhd G$ sedemikian sehingga $f$ berfaktor melalui
$(G/K_1)^n\to M$. Khususnya, $f$ hanya mengambil sejumlah hingga nilai;
kemulusan memberikan subgrup normal buka $K_2\lhd G$ sedemikian sehingga
nilai-nilai $f$ terletak dalam $M^{K_2}$. Ambil $K:=K_1\cap K_2$. Maka
\[ C^n(G,M)\simeq\varinjlim_{K\lhd G\;\text{buka}}C^n(G/K,M^K),\qquad
	\overline{C}^n(G,M)\simeq
	\varinjlim_{K\lhd G\;\text{buka}}\overline{C}^n(G/K,M^K). \]
Ketika $n$ bervariasi, ini memberikan isomorfisme kompleks.

% segment-id: o014.li2.chapter6.profinite-cohomology-a.p014
Seperti sebelumnya, definisikan submodul dari $C^n(G,M)$
\[ Z^n(G,M):=\Ker(d^n),\qquad B^n(G,M):=\Image(d^{n-1}), \]
dengan $B^0(G,M):=\{0\}$. Maka
$\Hm^n(C(G,M))=Z^n(G,M)/B^n(G,M)$. Definisikan dengan cara serupa versi
ternormalisasi $\overline{Z}^n(G,M)$ dan $\overline{B}^n(G,M)$.

% segment-id: o014.li2.chapter6.profinite-cohomology-a.t002
\begin{theorem}\label{prop:profinite-group-standard}
	Misalkan $G$ grup profinit. Untuk semua $n\in\Z_{\geq0}$ dan semua modul-$G$
	mulus $M$, terdapat isomorfisme kanonik
	\[ \Hm^n(C(G,M))\simeq\Hm^n(G,M)
	\simeq\Hm^n(\overline{C}(G,M)). \]
\end{theorem}
\begin{proof}
	Kohomologi bernilai dalam $\cate{Ab}$ atau $\Bbbk\dcate{Mod}$ jika
	gelanggang koefisien $\Bbbk$ turut ditetapkan. Kolimit terfilter bersifat
	eksak dalam kategori-kategori ini \cite[Teorema 6.2.2]{Li1}. Karena itu,
	\[ Z^n(G,M)=\varinjlim_KZ^n(G/K,M^K),\qquad
	B^n(G,M)=\varinjlim_KB^n(G/K,M^K), \]
	dan selanjutnya
	\begin{multline*}
		\Hm^n(C(G,M))
		=\dfrac{\varinjlim_KZ^n(G/K,M^K)}
		{\varinjlim_KB^n(G/K,M^K)}\simeq\\
		\varinjlim_K\frac{Z^n(G/K,M^K)}{B^n(G/K,M^K)}
		=\varinjlim_K\Hm^n(C(G/K,M^K)).
	\end{multline*}
	Terapkan Proposisi~\ref{prop:groupcoh-cochain} pada suku terakhir untuk
	memperoleh $\varinjlim_K\Hm^n(G/K,M^K)$, yakni $\Hm^n(G,M)$.

	Kasus $\overline{C}(G,M)$ serupa.
\end{proof}
